\begin{frame}
    \frametitle{How do I connect my instance?}

    You can interconnect your own MISP with the MN MISP instance in several ways ---
    pick what matches your \textbf{operational posture}.

    \vspace{0.5em}

    \begin{itemize}
        \item \textbf{Synchronisation} --- push and/or pull events between the two instances
            \begin{itemize}
                \item[] Initiated \textbf{by your instance}, or \textbf{by NCIA}
            \end{itemize}
        \item \textbf{Cross-correlation only} --- cache the remote data to spot overlaps,
              without ingesting it
        \item \textbf{Filtering} --- shape the flow so you only get what is relevant to you
    \end{itemize}

    \vspace{0.5em}

    \alertbox{info}{These modes are not exclusive}{Most members combine synchronisation
    with filtering to keep the flow under control.}
\end{frame}


\begin{frame}
    \frametitle{Two modes of synchronisation}

    Synchronisation can be driven from either side of the link.

    \vspace{0.5em}

    \begin{center}
        \begin{columns}[T]
            \begin{column}{0.49\textwidth}
                \card{\linewidth}{2.5em}
                    {{\Large\faArrowRight}\hspace{0.3em}\textbf{You initiate}}
                    {
                        \begin{itemize}
                            \item Your instance runs the \textbf{push / pull}
                            \item You control scheduling and direction
                            \item No inbound access to your instance required
                        \end{itemize}
                    }
            \end{column}
            \begin{column}{0.49\textwidth}
                \card{\linewidth}{2.5em}
                    {{\Large\faArrowLeft}\hspace{0.3em}\textbf{NCIA initiates}}
                    {
                        \begin{itemize}
                            \item NCIA pushes to you
                            \item \textbf{Real-time} on event publish
                            \item Your instance must be reachable
                        \end{itemize}
                    }
            \end{column}
        \end{columns}
    \end{center}
\end{frame}


\begin{frame}
    \frametitle{Mode 1 --- Your instance initiates (push / pull)}

    You reach out to the MN MISP instance and decide when and what to synchronise.

    \vspace{0.5em}

    \begin{enumerate}
        \item Ask NCIA for a \textbf{sync user} on the MN MISP instance
        \item Add \textbf{their instance} as a server on \textbf{yours}
        \item Decide your posture:
            \begin{itemize}
                \item[] \textbf{Ingest only} --- pull data down, keep quiet, \emph{or}
                \item[] \textbf{Ingest \& share back} --- also push your own contributions
            \end{itemize}
    \end{enumerate}

    \vspace{0.5em}

    \alertbox{info}{You stay in full control}{Scheduling, direction and what leaves your
    instance are all decided on your side.}
\end{frame}


\begin{frame}
    \frametitle{Mode 2 --- NCIA initiates the synchronisation}

    You give NCIA a \textbf{sync user of your own}, letting them push into your instance.

    \vspace{0.5em}

    \begin{itemize}
        \item Enables \textbf{real-time pushes} the moment an event is published by the
              MN MISP community
        \item Intel lands in your instance without waiting for your next pull cycle
        \item Requires your instance to be \textbf{reachable} --- from the internet or via an uplink
    \end{itemize}

    \vspace{0.5em}

    \alertbox{warning}{Reachability required}{Your instance must be exposed to NCIA ---
    over the internet or a dedicated uplink --- for inbound pushes to work.}
\end{frame}


\begin{frame}
    \frametitle{Cross-correlation only mode}

    Don't want to ingest the data at all? You can still benefit from it.

    \vspace{0.5em}

    \begin{itemize}
        \item Enable \textbf{caching} of the MN MISP instance
        \item Simply \textbf{see what NCIA holds} that may be relevant to you
        \item Correlations light up \textbf{locally}, without copying the full dataset into
              your own events
    \end{itemize}

    \vspace{0.5em}

    \alertbox{info}{Minimal footprint}{A lightweight way to know \emph{whether} the community
    has something on your indicators --- before deciding to pull it.}
\end{frame}


\begin{frame}
    \frametitle{Filtering --- staying in control of the flow}

    Worried about being overwhelmed with data? \textbf{Filter} what you take in.

    \vspace{0.5em}

    \begin{itemize}
        \item Constrain the pull by \textbf{tags, organisations, distribution}, \dots
        \item You \emph{can} filter out extremely old data --- though we \textbf{recommend against it}
    \end{itemize}

    \vspace{0.5em}

    \alertbox{warning}{On filtering out old data}{Operational use of 10+ year old events is
    very low --- \emph{but} for threat analysts, watching a threat actor evolve over a decade
    can be invaluable.}
\end{frame}
